\setcounter{chapter}{-1}
\chapter{Podstawowe pojęcia}
\label{chap:podstawowe-pojecia}
\poczatekwykladu{1}

\section{Odcinki i prostokąty}
\label{sec:terminologia}

\begin{definicja}[Odcinki]
Niech $a,b\in\R$ oraz $a\leq b$. Odcinek otwarty o końcach $a$ i $b$ to
\begin{equation}
  (a,b)=\{x\in\R:a<x<b\}.
  \label{eq:odcinek-otwarty}
\end{equation}
Analogicznie określamy odcinek domknięty i odcinki półotwarte:
\begin{align*}
  [a,b]&=\{x\in\R:a\leq x\leq b\},\\
  [a,b)&=\{x\in\R:a\leq x<b\},\\
  (a,b]&=\{x\in\R:a<x\leq b\}.
\end{align*}
W szczególności jednopunktowy zbiór $\{a\}$ zapisujemy również jako $[a,a]$.
\end{definicja}

\begin{definicja}[Prostokąt uogólniony]
Niech $d\in\N$ oraz niech $I_1,\ldots,I_d$ będą odcinkami w $\R$.
Prostokątem uogólnionym w $\R^d$ nazywamy zbiór
\begin{equation}
  P=I_1\times\cdots\times I_d
   =\{(x_1,\ldots,x_d)\in\R^d:
       x_i\in I_i\text{ dla każdego }i=1,\ldots,d\}.
  \label{eq:prostokat}
\end{equation}
\end{definicja}

\section{Prosta rozszerzona}
\label{sec:prosta-rozszerzona}

\begin{definicja}[Prosta rozszerzona]
Prostą rozszerzoną nazywamy zbiór
\begin{equation}
  \Rbar=[-\infty,+\infty]=\R\cup\{-\infty,+\infty\}.
  \label{eq:prosta-rozszerzona}
\end{equation}
Dla każdej liczby $a\in\R$ przyjmujemy
\[
  -\infty<a<+\infty.
\]
\end{definicja}

Dla $a\in\R$ obowiązują reguły
\begin{align*}
  a+(+\infty)&=(+\infty)+a=+\infty,\\
  a-(+\infty)&=-\infty.
\end{align*}
W teorii miary stosujemy także konwencję
\begin{equation}
  0\cdot(+\infty)=(+\infty)\cdot0=0.
  \label{eq:zero-nieskonczonosc}
\end{equation}

\begin{uwaga}
Równość \eqref{eq:zero-nieskonczonosc} jest przyjętą konwencją rachunku
na prostej rozszerzonej. Nie jest regułą obliczania granic: iloczyn typu
$0\cdot\infty$ w zadaniu o granicach pozostaje formą nieoznaczoną.
\end{uwaga}

\section{Granica górna i dolna ciągu zbiorów}
\label{sec:granice-zbiorow}

Niech $(A_n)_{n\geq1}$ będzie ciągiem podzbiorów ustalonego zbioru $X$.

\begin{definicja}[Granica górna]
Granicą górną ciągu $(A_n)$ nazywamy zbiór
\begin{equation}
  \limsup_{n\to\infty}A_n
  =\bigcap_{n=1}^{\infty}\bigcup_{k=n}^{\infty}A_k
  =\{x\in X:\forall n\geq1\ \exists k\geq n\quad x\in A_k\}.
  \label{eq:limsup}
\end{equation}
Innymi słowy, $x\in\limsup_{n\to\infty}A_n$ wtedy i tylko wtedy, gdy $x$
należy do nieskończenie wielu zbiorów $A_n$.
\end{definicja}

\begin{definicja}[Granica dolna]
Granicą dolną ciągu $(A_n)$ nazywamy zbiór
\begin{equation}
  \liminf_{n\to\infty}A_n
  =\bigcup_{n=1}^{\infty}\bigcap_{k=n}^{\infty}A_k
  =\{x\in X:\exists n\geq1\ \forall k\geq n\quad x\in A_k\}.
  \label{eq:liminf}
\end{equation}
Innymi słowy, $x\in\liminf_{n\to\infty}A_n$ wtedy i tylko wtedy, gdy $x$
należy do prawie wszystkich zbiorów $A_n$, czyli do wszystkich od pewnego
miejsca ciągu.
\end{definicja}

\begin{fakt}
Dla każdego ciągu zbiorów $(A_n)$ zachodzi
\begin{equation}
  \liminf_{n\to\infty}A_n\subseteq\limsup_{n\to\infty}A_n.
  \label{eq:inkluzja-granic}
\end{equation}
\end{fakt}

\begin{proof}
Jeżeli punkt należy do wszystkich zbiorów od pewnego miejsca ciągu,
to należy do nieskończenie wielu zbiorów tego ciągu.
\end{proof}

\begin{definicja}[Zbieżność ciągu zbiorów]
Jeżeli
\[
  \liminf_{n\to\infty}A_n=\limsup_{n\to\infty}A_n,
\]
to mówimy, że ciąg $(A_n)$ jest zbieżny. Wspólną wartość tych granic
oznaczamy przez $\lim_{n\to\infty}A_n$.
\end{definicja}

\section{Przykłady granic ciągów zbiorów}
\label{sec:przyklady-granic}

\begin{przyklad}[Ciąg naprzemienny]
Niech
\[
  A_n=[0,1+(-1)^n],\qquad n\geq1.
\]
Wówczas
\[
  A_{2k-1}=\{0\},\qquad A_{2k}=[0,2],\qquad k\geq1.
\]
Punkt $0$ należy do każdego zbioru $A_n$. Każdy punkt $x\in(0,2]$ należy
do wszystkich zbiorów o indeksach parzystych, a nie należy do żadnego
zbioru o indeksie nieparzystym. Punkty spoza $[0,2]$ nie należą do
żadnego z tych zbiorów. Zatem, zgodnie z \eqref{eq:limsup} i \eqref{eq:liminf},
\begin{equation}
  \liminf_{n\to\infty}A_n=\{0\},\qquad
  \limsup_{n\to\infty}A_n=[0,2].
  \label{eq:przyklad-naprzemienny}
\end{equation}
Granice są różne, więc ciąg $(A_n)$ nie jest zbieżny.
\end{przyklad}

\begin{przyklad}[Początkowe odcinki zbioru liczb naturalnych]
Przyjmujemy $\N=\{1,2,3,\ldots\}$ i rozważamy ciąg
\[
  A_n=\{1,2,\ldots,n\},\qquad n\geq1.
\]
Dla każdego $m\in\N$ mamy $m\in A_n$ dla wszystkich $n\geq m$.
Każda liczba naturalna należy więc do prawie wszystkich zbiorów $A_n$.
Punkty spoza $\N$ nie należą do żadnego z nich. Stąd
\begin{equation}
  \liminf_{n\to\infty}A_n=\N=\limsup_{n\to\infty}A_n,
  \qquad \lim_{n\to\infty}A_n=\N.
  \label{eq:przyklad-naturalne}
\end{equation}
\end{przyklad}
